\documentclass[11pt]{article}

% ===== 基本宏包：尽量保持简单 =====
\usepackage[margin=1in]{geometry}
\usepackage{graphicx}
\usepackage{float}
\usepackage{chngcntr}
\usepackage{booktabs}
\usepackage{multirow}
\usepackage{longtable}
\usepackage{array}
\usepackage{calc}
\usepackage{amsmath, amssymb}
\setlength{\LTcapwidth}{\textwidth}
\usepackage{xurl}
\usepackage{hyperref}
\usepackage[numbers,sort&compress]{natbib}
\usepackage{setspace}
\usepackage{fontspec}
\usepackage{xeCJK}
\providecommand{\tightlist}{\setlength{\itemsep}{0pt}\setlength{\parskip}{0pt}}

% ===== 图表编号风格：全文连续编号，使用 Figure / Table =====
\renewcommand{\figurename}{Figure}
\renewcommand{\tablename}{Table}

% ===== 字体与换行策略 =====
\defaultfontfeatures{Ligatures=TeX,Scale=MatchLowercase}
\IfFontExistsTF{Noto Serif}{\setmainfont{Noto Serif}}{\IfFontExistsTF{DejaVu Serif}{\setmainfont{DejaVu Serif}}{\setmainfont{Latin Modern Roman}}}
\IfFontExistsTF{Noto Sans}{\setsansfont{Noto Sans}}{\IfFontExistsTF{DejaVu Sans}{\setsansfont{DejaVu Sans}}{\setsansfont{Latin Modern Sans}}}
\IfFontExistsTF{Noto Sans Mono}{\setmonofont{Noto Sans Mono}}{\IfFontExistsTF{DejaVu Sans Mono}{\setmonofont{DejaVu Sans Mono}}{\setmonofont{Latin Modern Mono}}}
\IfFontExistsTF{Noto Serif CJK SC}{\setCJKmainfont[AutoFakeBold=3,AutoFakeSlant=.2]{Noto Serif CJK SC}}{\IfFontExistsTF{Source Han Serif SC}{\setCJKmainfont[AutoFakeBold=3,AutoFakeSlant=.2]{Source Han Serif SC}}{\IfFontExistsTF{AR PL UMing CN}{\setCJKmainfont[AutoFakeBold=3,AutoFakeSlant=.2]{AR PL UMing CN}}{\setCJKmainfont[AutoFakeBold=3,AutoFakeSlant=.2]{WenQuanYi Zen Hei}}}}
\IfFontExistsTF{Noto Sans CJK SC}{\setCJKsansfont[AutoFakeBold=3,AutoFakeSlant=.2]{Noto Sans CJK SC}}{\IfFontExistsTF{Source Han Sans SC}{\setCJKsansfont[AutoFakeBold=3,AutoFakeSlant=.2]{Source Han Sans SC}}{\IfFontExistsTF{Droid Sans Fallback}{\setCJKsansfont[AutoFakeBold=3,AutoFakeSlant=.2]{Droid Sans Fallback}}{\setCJKsansfont[AutoFakeBold=3,AutoFakeSlant=.2]{AR PL UMing CN}}}}
\IfFontExistsTF{Noto Sans Mono CJK SC}{\setCJKmonofont[AutoFakeBold=3]{Noto Sans Mono CJK SC}}{\IfFontExistsTF{Droid Sans Fallback}{\setCJKmonofont[AutoFakeBold=3]{Droid Sans Fallback}}{\setCJKmonofont[AutoFakeBold=3]{AR PL UMing CN}}}
\setlength{\emergencystretch}{6em}
\sloppy

% 可选：让行距略宽一些，便于审稿阅读
\onehalfspacing

\title{\textbf{Decathlon VOC Analyzer：面向商品图文证据与用户评论对齐的证据驱动多模态 VOC 分析系统}}

\author{{\small
\begin{tabular}{@{}ccc@{}}
\multicolumn{3}{c}{\parbox[t]{0.6600\textwidth}{\centering \textbf{Severin Ye}\textsuperscript{1} \\ \textit{School of Computer Science and Engineering} \\ \textit{Kyungpook National University} \\ Daegu, Republic of Korea \\ {\ttfamily\small 6severin9\allowbreak@gmail.\allowbreak{}com}}} \\
\multicolumn{3}{c}{} \\[0.9em]
\parbox[t]{0.3067\textwidth}{\centering \textbf{Dokeun Lee}\textsuperscript{2} \\ \textit{Department of Physics} \\ \textit{Kyungpook National University} \\ Daegu, Republic of Korea \\ {\ttfamily\small nsa08008\allowbreak@naver.\allowbreak{}com}} & \parbox[t]{0.3067\textwidth}{\centering \textbf{Wushuang Liu}\textsuperscript{2} \\ \textit{Department of Child Studies} \\ \textit{Kyungpook National University} \\ Daegu, Republic of Korea \\ {\ttfamily\small 824309203q\allowbreak@gmail.\allowbreak{}com}} & \parbox[t]{0.3067\textwidth}{\centering \textbf{Seowan Jung}\textsuperscript{2} \\ \textit{School of Computer Science and Engineering} \\ \textit{Kyungpook National University} \\ Daegu, Republic of Korea \\ {\ttfamily\small swan041014\allowbreak@gmail.\allowbreak{}com}} \\[0.9em]
\multicolumn{3}{c}{\makebox[0.9200\textwidth][c]{\parbox[t]{0.3067\textwidth}{\centering \textbf{HyunJun Jung}\textsuperscript{2} \\ \textit{School of Computer Science and Engineering} \\ \textit{Kyungpook National University} \\ Daegu, Republic of Korea \\ {\ttfamily\small beanbeansummon\allowbreak@naver.\allowbreak{}com}}\hspace{2\tabcolsep}\parbox[t]{0.3067\textwidth}{\centering \textbf{Hye Jeon}\textsuperscript{2} \\ \textit{School of Computer Science and Engineering} \\ \textit{Kyungpook National University} \\ Daegu, Republic of Korea \\ {\ttfamily\small prajna0426\allowbreak@naver.\allowbreak{}com}}}} \\
\end{tabular}
}}

\date{} % Nature 系列投稿一般不需要显示日期

\begin{document}
\maketitle
\nocite{lewis2020retrievalAugmentedGeneration,gao2023rarr,radford2021clip,liu2024structuredOutputBenchmark,geng2023grammarConstrainedDecoding,yao2023react,nvidia2024multimodalRagIntro,nvidia2024aiReadyKnowledgeSystems,faysse2024colpali,pontiki2014semevalAspect,qdrant2024pdfRetrievalAtScale,qdrant2024multiVectorCourse,manning2008introductionIr}

% =========================
% Abstract
% =========================
\begin{abstract}
电商用户之声（VOC）分析需同时理解评论主观体验与商品页客观证据。本文提出 Decathlon VOC Analyzer，一个证据驱动的多模态 VOC 分析系统。该方法将商品信息标准化为证据包，对评论进行方面级抽取，将每个方面规划为具有明确检索路线的问题，执行两阶段多模态召回与重排，最终生成带 claim 级证据归因的结构化报告。系统保留关键中间表示以支持复现和误差分析。本文贡献包括：提出``评论方面建模---问题规划---多模态检索---证据约束生成''框架；实现结构化中间表示的可复现系统；通过 166 项自动化测试验证系统稳定性。本文定位为系统方法论文，不声称已完成大规模人工标注 benchmark。

\subsection{关键词}

多模态 VOC 分析；证据约束生成；图文检索；商品评论挖掘；主张归因
\end{abstract}

% =========================
% Ⅰ. 引言
% =========================
\section{Ⅰ. 引言}
用户评论分析是电商智能运营中的基础任务。传统方法通常围绕情感分类、主题聚类或方面级情感分析展开，能够概括用户对商品属性的正负态度，却较少回答这些评论观点是否能够被商品本体证据解释、支持或反驳。对于商品运营和产品改进而言，单纯知道``用户抱怨容量小''并不足够，分析者还需要知道商品页面是否明确描述了尺寸、图片是否展示了空间结构、该问题是商品缺陷还是用户预期偏差。

大语言模型为商品评论总结提供了新的可能，但一步式总结存在明显风险。如果系统直接将商品文案、图片描述和评论整体输入模型，模型可能生成流畅但难以审查的结论。使用者很难判断某条建议来自哪些评论、哪些商品文本或哪张图片，也难以定位错误发生在评论理解、检索、证据解释还是最终生成阶段。因此，商品 VOC 分析不应被简化为黑盒摘要任务，而应被建模为一条从评论需求到商品证据、再到结构化报告的可追踪链路。

本文研究的问题是：如何在电商商品场景中，将用户评论中的主观体验显式对齐到商品文案和商品图像中的可核验证据，并生成可审查、可复现、可评估的 VOC 洞察。为此，本文提出 Decathlon VOC Analyzer。系统以商品为中心构建标准化证据包，将评论抽取为方面级对象，再通过问题规划层把每个方面转化为文本、图像或图文联合检索任务，最后通过多模态召回、重排、证据聚合和 claim 级归因生成结构化报告。

% =========================
% Ⅱ. 方法章节
% =========================
\section{Ⅱ. 方法章节}
本文将商品 VOC 分析组织为一个证据驱动的方法章节。核心目标不是直接生成评论摘要，而是把评论中的主观反馈转化为可检索、可验证、可归因的图文证据链路，并在统一结构下完成问题规划、证据检索、实验说明与图示展示。

\subsection{总体架构}

Decathlon VOC Analyzer 将商品 VOC 分析形式化为证据驱动的多阶段流程。系统首先把商品页面中的文本、图像和评论组织为统一证据空间；随后在评论侧抽取方面级用户反馈；再通过问题规划将主观反馈转化为可执行的证据查询；最后在商品文本与商品图像之间执行多模态检索、重排和证据约束生成 \citep{lewis2020retrievalAugmentedGeneration,gao2023rarr,radford2021clip,nvidia2024multimodalRagIntro,nvidia2024aiReadyKnowledgeSystems,faysse2024colpali}。与一步式大模型摘要不同，该框架将``用户表达了什么''``商品证据能否支持''``最终报告如何归因''拆分为可观察的中间环节。

系统包含四个逻辑层。商品证据层负责构建商品文本、图片和局部视觉区域的结构化表示。VOC 需求层负责将评论转化为方面、观点、情感和场景。证据对齐层负责把方面改写为具有明确模态路线的问题，并在商品证据空间中检索和重排候选。报告归因层负责生成结构化洞察，并将报告主张映射回评论证据与商品证据。

\subsection{公式}

为统一评论方面 \(a\) 与候选证据 \(e\) 的排序，本文将文本相关性、图像相关性与跨模态一致性组合为一个证据评分：

\[
S(a, e) = \lambda_t R_t(a, e) + \lambda_v R_v(a, e) + \lambda_c C(a, e)
\]

其中，\(R_t(a, e)\) 表示文本路线的相关性得分，\(R_v(a, e)\) 表示图像路线的相关性得分，\(C(a, e)\) 表示评论方面与商品证据之间的跨模态一致性约束，\(\lambda_t\)、\(\lambda_v\) 和 \(\lambda_c\) 为可调权重。该表达式的作用不是替代最终的生成式分析，而是在检索与重排阶段提供统一的证据选择准则。

\subsection{实验说明}

在商品证据表示阶段，系统将商品标题、描述、类目、规格、图片和用户评论表示为具有稳定标识的证据对象。文本证据保留来源字段，图像证据同时保留整图和规则化局部区域，评论证据保留评分、语言提示和原文片段，从而为后续检索、归因和评估建立统一对象边界。

在评论建模阶段，系统将自然语言反馈转化为方面级 VOC 单元 \citep{pontiki2014semevalAspect}。每个方面单元包含被评价属性、情感倾向、观点描述、原文证据片段、使用场景和置信度。方面级表示既支持统计聚合，也支持围绕具体属性生成证据查询。为避免有限评论预算被高星短评占据，系统在评论抽取前引入星级感知抽样，并支持结构化 LLM 路径与启发式路径，以区分模型能力和流程完整性。

在问题规划阶段，系统不直接用原始评论检索商品证据，而是将每个方面转化为具有明确证据意图的问题。这些问题覆盖文本支持、视觉确认和跨模态解释三类需求：文本支持问题关注商品文案或规格是否支持评论观点，视觉确认问题关注商品图片是否呈现相关结构或外观，跨模态解释问题则判断某一反馈更接近真实产品问题、页面呈现不足还是用户预期偏差。

在检索与生成阶段，系统采用两阶段检索 \citep{qdrant2024pdfRetrievalAtScale,qdrant2024multiVectorCourse}。第一阶段使用向量表示从商品文本和商品图像中粗召回候选证据，第二阶段使用重排模型对候选集合进行精排。文本路线主要用于验证商品名称、类目、描述、规格和显式承诺，图像路线主要用于验证结构、外观、颜色和局部细节。最终报告在方面聚合基础上生成优势、缺点、争议点、适用场景、证据缺口和改进建议，并通过 claim 级归因标记主张与评论、商品文本或商品图像之间的支持关系。

系统实现采用结构化对象约束多阶段流程 \citep{liu2024structuredOutputBenchmark,geng2023grammarConstrainedDecoding,yao2023react}，模型调用通过兼容式语言模型网关完成，检索层支持本地索引和向量数据库两类后端。为了区分正式实验与开发验证，系统提供运行策略控制：当完整模型链路可用时使用真实向量编码、图像编码和重排模型；当外部能力不可用且策略允许时，退回启发式路径以保证流程可运行。该设计避免将模型能力不足和流程设计失败混淆在一起。

本文当前验证目标是确认系统是否能够稳定完成证据驱动 VOC 分析，而不是报告冻结大规模 benchmark 的最终性能。结果表明，当前实现能够依次完成商品证据标准化、评论方面建模、问题规划、多模态召回、候选重排、方面聚合、报告生成和证据归因。一次完整分析不仅包含最终自然语言报告，还包含评论方面、问题意图、检索问题、召回证据、检索质量、运行配置、方面聚合、过程追踪和报告归因等对象，因而支持人工审查和错误定位。

当存在问题级相关证据标注时，系统可计算 Recall@1、Recall@3、Recall@5、MRR、NDCG@3 和 NDCG@5 \citep{manning2008introductionIr}。对于结构化报告，系统统计主张是否被评论、商品文本或商品图像支持，包括 claim support rate、claim grounded rate、citation precision、citation contradiction rate 和不同模态路线的贡献比例。当前实现包含 166 项自动化测试且全部通过，说明系统主要接口、结构化中间表示和工作流断言在当前实现中保持一致。

\subsection{图 1}

\begin{figure}[H]
\centering
\includegraphics[width=0.92\textwidth]{/home/severin/Codelib/Decathlon\_VOC\_Analyzer/论文\_md形式/图片/图1.png}
\caption{Decathlon VOC Analyzer 的总体框架。系统以商品证据包和评论方面为中间表示，通过问题规划连接主观 VOC 信号与多模态商品证据，并在报告阶段进行 claim 级归因。}
\end{figure}

\subsection{图 2}

\begin{figure}[H]
\centering
\includegraphics[width=0.92\textwidth]{/home/severin/Codelib/Decathlon\_VOC\_Analyzer/论文\_md形式/图片/图2.png}
\caption{评论方面到证据查询的规划机制。方面对象并不直接进入检索，而是先被分解为具有不同证据意图的问题，再分别触发文本路线、图像路线或图文联合路线。}
\end{figure}

% =========================
% Ⅲ. 结论
% =========================
\section{Ⅲ. 结论}
本文提出 Decathlon VOC Analyzer，一个面向商品图文证据与用户评论对齐的证据驱动多模态 VOC 分析系统。系统将原始商品数据标准化为可追踪证据包，将用户评论抽取为方面级 VOC 对象，再通过问题规划、多模态召回、重排、报告生成和 claim 级归因生成结构化商品分析结果。

本文的核心结论是：商品 VOC 分析不应被简化为一步式评论摘要，而应被组织为``评论方面建模---证据问题规划---商品图文检索---证据约束生成''的多阶段流程。当前实现已经形成可运行的研究原型，并在结构化中间表示、证据归因、回放机制和评估接口上保持一致。测试结果显示当前实现通过 166 项自动化测试，说明系统主链和关键模块在工程层面具备稳定性。

未来工作将重点补充多品类人工标注 benchmark，系统比较不同检索和重排策略，引入更细粒度视觉 grounding，并将人工反馈与 replay 机制进一步闭环。

% =========================
% References
% =========================
\bibliographystyle{unsrtnat}
\bibliography{ref}

\end{document}
